%%%%%%%%%%%%%%%%%%%%%%%%%%%%%%%%%%%%%%%%%%%%%%%%%%%%%%%%%%%%%%%%%%%%%%
% Lớp 10 | Thống kê | Mô tả và biểu diễn dữ liệu trên các bảng và biểu đồ
% Dạng 0D6?2-2: Đọc và phân tích thông tin trên biểu đồ
% Số câu: 13  (sinh tự động bởi tach_cau.py; từ nay thầy sửa trực tiếp file này)
%%%%%%%%%%%%%%%%%%%%%%%%%%%%%%%%%%%%%%%%%%%%%%%%%%%%%%%%%%%%%%%%%%%%%%

%%%%%%%%%%%%%%%% TRẮC NGHIỆM 4 PHƯƠNG ÁN %%%%%%%%%%%%%%%%

% nguồn: MaHoaNH_TracNghiem_Lop10
\begin{ex}%[0D6K2-2]%[ma:9dfe9f91]
	Trên nửa đường tròn đơn vị cho góc $\alpha$  sao cho $\sin\alpha=\dfrac{2}{3} $ và $\cos\alpha<0$. Tính $\tan\alpha$.
	\choice 
		 {\True $-\dfrac{2\sqrt{5}}{5}$}
		 {$\dfrac{2\sqrt{5}}{5}$}
		 {$-\dfrac{2}{5}$}
		{$1$} 
	 \loigiai{
		Ta có $\cos\alpha=-\sqrt{1-\sin^2\alpha}=-\dfrac{\sqrt{5}}{3}$ suy ra $\tan\alpha=-\dfrac{2\sqrt{5}}{5}$. 
	} 
\end{ex}

% nguồn: Ok_TracNghiemLop11_Lan1
\begin{ex}%[0D6K2-2]%[ma:3e466769]
	 Cho góc $\alpha $ thỏa mãn $\tan \alpha =2.$ Tính $P=\dfrac{3\sin \alpha-2\cos \alpha}{5\cos \alpha+7\sin \alpha}.$
	\choice 
		 {$P=-\dfrac{4}{9}$}
		{$P=\dfrac{4}{9}$} 
		 {$P=-\dfrac{4}{19}$}
		{\True $P=\dfrac{4}{19}$} 
	\loigiai{Chia cả tử và mẫu của $P$ cho $\cos  \alpha $ ta được $P=\dfrac{3\tan \alpha-2}{5+7\tan \alpha}=\dfrac{3\cdot 2-2}{5+7\cdot 2}=\dfrac{4}{19}.$
	 }
\end{ex}

% nguồn: Ok_TracNghiemLop11_Lan1
\begin{ex}%[0D6B2-2]%[ma:9261d427]
	 Cho góc $\alpha $ thỏa mãn $\sin \alpha =\dfrac{12}{13}$ và $\dfrac{\pi}{2}<\alpha <\pi $. Tính $\cos \alpha.$
	 \choice
		{$\cos  \alpha =\dfrac{1}{13}$}
		{$\cos \alpha  =\dfrac{5}{13}$}
		 {\True $\cos \alpha =-\dfrac{5}{13}$}
		{$\cos \alpha =-\dfrac{1}{13}$} 
	\loigiai{Ta có $\heva{&  \cos \alpha =\pm \sqrt{{1-{\sin}^2\alpha}}=\pm \dfrac{5}{13} \\
		& \dfrac{\pi}{2}<\alpha <\pi}\Rightarrow\cos \alpha  =-\dfrac{5}{13}.$
	 }
\end{ex}

% nguồn: Ok_TracNghiemLop11_Lan1
\begin{ex}%[0D6B2-2]%[ma:ad6eed1a]
	Cho góc $\alpha  $ thỏa mãn $\tan \alpha =-\dfrac{4}{3}$ và $\dfrac{2017\pi}{2}<\alpha <\dfrac{2019\pi}{2}$. Tính $\sin \alpha.$
	\choice 
		{$\sin \alpha =-\dfrac{3}{5}$} 
		{$\sin \alpha =\dfrac{3}{5}$} 
		{$\sin \alpha =-\dfrac{4}{5}$} 
		{\True $\sin  \alpha =\dfrac{4}{5}$}
	\loigiai{Ta có $\heva{& 1+\tan  ^2\alpha =\dfrac{1}{\cos^2\alpha} \\
		& \dfrac{2017\pi}{2}<\alpha  <\dfrac{2019\pi}{2}}\overset{\longleftrightarrow}{}\heva{& 1+\left(-\dfrac{4}{3}\right)^2=\dfrac{1}{\cos^2\alpha} \\
		& \dfrac{\pi}{2}+504 \cdot 2\pi <\alpha <\dfrac{3\pi}{2}+504 \cdot  2\pi}$
		$\Rightarrow\cos \alpha =-\dfrac{3}{5}$. Mà $\tan  \alpha =\dfrac{\sin \alpha}{\cos \alpha}=-\dfrac{4}{3}=\dfrac{\sin \alpha}{-\dfrac{3}{5}}\Rightarrow\sin \alpha =\dfrac{4}{5}$.
	 }
\end{ex}

% nguồn: Ok_TracNghiemLop11_Lan1
\begin{ex}%[0D6B2-2]%[ma:dfbabf0d]
	Cho góc $\alpha $ thỏa mãn $\cos \alpha  =-\dfrac{12}{13}$ và $\dfrac{\pi}{2}<\alpha <\pi.$ Tính $\tan \alpha.$
	 \choice
		 {$\tan \alpha =-\dfrac{12}{5}$}
		 {$\tan \alpha =\dfrac{5}{12}$}
		{\True $\tan \alpha  =-\dfrac{5}{12}$}
		{$\tan \alpha =\dfrac{12}{5}$} 
	\loigiai{Ta có $\heva{& \sin \alpha =\pm \sqrt{{1-{\cos}^2\alpha}}=\pm \dfrac{5}{13}  \\
		&  \dfrac{\pi}{2}<\alpha <\pi}\Rightarrow\alpha =\dfrac{5}{13}\Rightarrow\tan \alpha =\dfrac{\sin \alpha}{\cos \alpha}=-\dfrac{5}{12}.$
	} 
\end{ex}

% nguồn: Ok_TracNghiemLop11_Lan1
\begin{ex}%[0D6K2-2]%[ma:5bbd82da]
	Cho góc $\alpha $ thỏa mãn $\cot\alpha =\dfrac{1}{3}$. Tính  $P=\dfrac{3\sin \alpha+4\cos \alpha}{2\sin \alpha-5\cos \alpha}.$
	\choice 
		 {$P=-\dfrac{15}{13}$}
		{$P=\dfrac{15}{13}$} 
		{$P=-13$} 
		 {\True $P=13$}
	\loigiai{ Chia cả tử  và mẫu của $P$ cho $\sin \alpha $ ta được $P=\dfrac{3+4\cot \alpha}{2-5\cot \alpha}=\dfrac{3+4\cdot \dfrac{1}{3}}{2-5\cdot \dfrac{1}{3}}=13$.
	} 
\end{ex}

% nguồn: Ok_TracNghiemLop11_Lan1
\begin{ex}%[0D6B2-2]%[ma:60407214]
	Cho góc $\alpha  $ thỏa mãn $\cos \alpha =-\dfrac{\sqrt{5}}{3}$ và $\pi <\alpha <\dfrac{3\pi}{2}$. Tính $\tan \alpha.$
	\choice 
		{$\tan \alpha  =-\dfrac{3}{\sqrt{5}}$}
		 {\True $\tan \alpha =\dfrac{2}{\sqrt{5}}$}
		{$\tan  \alpha =-\dfrac{4}{\sqrt{5}}$}
		 {$\tan \alpha =-\dfrac{2}{\sqrt{5}}$}
	\loigiai{ Ta có  $\heva{& \sin \alpha =\pm \sqrt{{1-{\cos}^2\alpha}}=\pm \dfrac{2}{3} \\
		& \pi <\alpha <\dfrac{3\pi}{2}}\Rightarrow\alpha =-\dfrac{2}{3}\Rightarrow\tan \alpha =\dfrac{\sin\alpha}{\cos \alpha}=\dfrac{2}{\sqrt{5}}.$ 
	 }
\end{ex}

% nguồn: Ok_TracNghiemLop11_Lan1
\begin{ex}%[0D6B2-2]%[ma:e274f2b8]
	Cho góc $\alpha $ thỏa mãn  $\sin \alpha =\dfrac{3}{5}$ và $\dfrac{\pi}{2}<\alpha <\pi $. Tính $P=\dfrac{\tan \alpha}{1+{\tan}^2\alpha}.$
	 \choice
		{$P=-3$} 
		{$P=\dfrac{3}{7}$} 
		 {$P=\dfrac{12}{25}$}
		{\True  $P=-\dfrac{12}{25}$}
	\loigiai{Ta có $\heva{& \cos  \alpha =\pm \sqrt{{1-{\sin}^2\alpha}}=\pm \dfrac{4}{5} \\
		& \dfrac{\pi}{2}<\alpha <\pi}\Rightarrow\cos \alpha =-\dfrac{4}{5}\Rightarrow\tan \alpha  =-\dfrac{3}{4}$.\\
		Thay  $\tan \alpha =-\dfrac{3}{4}$ vào $P$, ta được $P=-\dfrac{12}{25}$.
	} 
\end{ex}


%%%%%%%%%%%%%%%% ĐÚNG/SAI %%%%%%%%%%%%%%%%

% nguồn: MaHoaNH_DungSai_Lop10
\begin{ex}%[0D6H2-2]%[ma:c53bf5a4]
	\immini{ 
	Biểu đồ hình quạt tròn ở hình bên biểu diễn kết quả thống kê (tính theo tỉ số phần trăm) chọn môn thể thao ưa thích nhất trong bốn môn: Bóng đá, Cầu lông, Bóng bàn, Bóng rổ của học sinh khối 11 ở trường THPT X. Mỗi học sinh chỉ được chọn một môn thể thao khi được hỏi ý kiến. 
	 Biết rằng tổng số học sinh khối 11 là 500.
	Các khẳng  định sau đây đúng hay sai?
	 }
	{ 
	\begin{tikzpicture}[scale=0.85,font=\footnotesize,line join=round,line cap=round,>=stealth]
			\def\R{3} 
			\draw[double](0,0)  circle (\R);
			 \def\X{10/cyan/Bóng bàn,20/red/Bóng rổ,30/green/Bóng đá,40/teal/Cầu lông}
			\pgfmathsetmacro{\a}{0} 
			\foreach \x/\mau/\yeuto [count=\i] in  \X{
				\pgfmathsetmacro{\g}{\x/100*360} 
				\pgfmathsetmacro{\b}{\a-\g} 
				 \fill[\mau,opacity=0.5,draw=black](0,0)--({\b+\g}:\R) arc ({\b+\g}:\b:\R)--cycle;
				 \path ({\b+\g}:0.65*\R) arc ({\b+\g}:\b:0.65*\R) node[pos=0.5]{\fontsize{8pt}{6pt}\selectfont{\yeuto}}
				({\b+\g}:0.65*\R) arc ({\b+\g}:\b:0.65*\R) node[pos=0.5,below=2pt]{\fontsize{8pt}{6pt}\selectfont{\x $\%$}}; 
				\global \let \a  = \b;
			 }
		\end{tikzpicture}
	 }
	\choiceTF 
		{\True Cầu lông là môn thể thao  ưa thích nhất}
		{Số học  sinh ưa thích Bóng bàn là $40$ học sinh}
		{\True Số học sinh ưa thích Bóng rổ ít hơn số học sinh ưa thích Cầu  lông là 100 học sinh}
		{Số  học sinh ưa thích Bóng đá hoặc Bóng rổ là $200$ học sinh}
	 \loigiai{
		\begin{itemchoice}
			\itemch  Số học sinh ưa thích Bóng bàn là $10\% \cdot 500 =50$ học sinh.
			\itemch Số học sinh ưa thích  Bóng rổ ít hơn số học sinh ưa thích Cầu lông là $(40\%-20\%) \cdot 500 =100$ học sinh.
			\itemch Số học sinh ưa thích Bóng đá hoặc Bóng rổ là $(30\%+20\%) \cdot 500 =250$  học sinh.
		\end{itemchoice}
	 }
\end{ex}

% nguồn: MaHoaNH_DungSai_Lop10
\begin{ex}%[0D6H2-2]%[ma:5d92e593]
	Điểm thi đánh giá tư duy của bạn Đức  được biểu diễn bằng biểu đồ
	\begin{center}
		\begin{tikzpicture}[y=0.5cm,x=0.9cm]
				\foreach \i in {0,10,20,...,100} 
				\draw (-0.4,{\i/10}) node{\textcolor{black}{\fontsize{9}{9}\selectfont\bfseries \i}}; 
				\draw  [->] (0,0)--(11,0) node[below]{\text{Lần}};
				\draw [->] (0,0)--(0,12)  node[left]{\text{Điểm}};
				\foreach \i in  {1,...,10}
				\draw[color=gray!50,line width=1.5pt] (0,\i)--(11,\i); 
				\foreach \i in {1,3,5,7,9}  \draw (\i,-2pt)--(\i,2pt);
				 \draw
					(1,0) node[below]{\text{Lần 1}} 
					(3,0) node[below]{\text{Lần 2}} 
					(5,0)  node[below]{\text{Lần 3}}
					(7,0) node[below]{\text{Lần 4}} 
					 (9,0) node[below]{\text{Lần 5}};
				\draw[thick] (1,6)--(3,6.8)--(5,8.1)--(7,8.8)--(9,8.4); 
				\draw[fill=black] 
					(1,6) circle (1.2pt) node[above]{$60$} 
					 (3,6.8) circle (1.2pt) node[above]{$68$}
					(5,8.1) circle (1.2pt) node[above]{$81$} 
					(7,8.8)  circle (1.2pt) node[above]{$88$}
					(9,8.4)  circle (1.2pt) node[above]{$84$};
			\end{tikzpicture}
	\end{center}
	Các  khẳng định sau đây đúng hay sai?
	 \choiceTF
		{\True Điểm cao  nhất là $88$}
		{Tỉ lệ giữa điểm cao nhất và thấp nhất  lớn hơn $150\%$}
		{\True Hiệu số điểm hai lần thi liên tiếp lớn nhất là  $8$ điểm}
		{\True So với lần thi thứ 4, điểm thi lần thứ 5 giảm ít hơn $5\%$} 
	 \loigiai{
		\begin{itemchoice}
			\itemch Điểm cao nhất  là $88$.
			\itemch Tỉ lệ  giữa điểm cao nhất và thấp nhất là $\dfrac{88}{60} \approx 146{,}7\%$.
			\itemch Hiệu  số điểm hai lần thi liên tiếp lớn nhất là $68-60=8$ điểm.
			\itemch So với lần thi thứ 4, điểm thi lần thứ 5 giảm $\dfrac{4}{88}  \approx 4{,}5\%$.
		\end{itemchoice}
	} 
\end{ex}


%%%%%%%%%%%%%%%% TRẢ LỜI NGẮN %%%%%%%%%%%%%%%%

% nguồn: MaHoaNH_TraLoiNgan_Lop10
\begin{bt}%[0D6H2-2]%[ma:1d6e3d51]
	Trong giờ học Toán lớp $7$, giáo viên giao nhiệm vụ cho mỗi nhóm đo các góc của một tam giác. Kết  quả được ghi lại trong bảng sau
	\begin{center}
		\begin{tabular}{|c|c|c|c|c|}
			\hline 
			Nhóm & Nhóm $1$ & Nhóm $2$ &  Nhóm $3$ & Nhóm $4$\\
			 \hline
			Góc thứ nhất & $35^\circ$ &  $61^\circ$ & $33^\circ$ & $100^\circ$\\
			\hline 
			Góc thứ hai & $77^\circ$ & $74^\circ$ & $102^\circ$  & $37^\circ$\\
			 \hline
			Góc thứ ba & $68^\circ$ & $45^\circ$ &  $47^\circ$ & $43^\circ$\\
			\hline 
		\end{tabular}
	\end{center}
	Trong bảng trên có nhóm  ghi kết quả sai. Hãy tìm nhóm nào sai?
	\shortans{3} 
	\loigiai{ 
		Do tổng ba  góc trong một tam giác bằng $180^\circ$ nên ta lập bảng sau
		\begin{center}
			\begin{tabular}{|c|c|c|c|c|}
				 \hline
				Nhóm & Nhóm $1$ & Nhóm $2$ &  Nhóm $3$ & Nhóm $4$\\
				\hline 
				Góc thứ nhất & $35^\circ$ & $61^\circ$ & $33^\circ$  & $100^\circ$\\
				\hline 
				Góc thứ hai & $77^\circ$ & $74^\circ$ &  $102^\circ$ & $37^\circ$\\
				 \hline
				Góc thứ ba & $68^\circ$ & $45^\circ$  & $47^\circ$ & $43^\circ$\\
				\hline 
				Tổng ba  góc & $180^\circ$ & $180^\circ$ & $182^\circ$ & $180^\circ$\\
				 \hline
			\end{tabular}
		\end{center}
		Vậy nhóm $3$  đo sai vì tổng số đo của ba góc khác $180^\circ$.
	} 
\end{bt}


%%%%%%%%%%%%%%%% TỰ LUẬN %%%%%%%%%%%%%%%%

% nguồn: Ok_TuanLuan_Toan11_Lan1
\begin{bt}%[0D6K2-2]%[ma:28368f29]
	Cho $\cot\alpha =3$. Tính giá trị biểu thức $M=\dfrac{2\sin \alpha-3\cos\alpha}{5\sin^3 \alpha+\cos^3  \alpha}$.
	 \loigiai{
		Ta có 
		\begin{align*}
			M&=\dfrac{2\sin  \alpha-3\cos\alpha}{5\sin^3 \alpha+\cos^3 \alpha}\\
			 &=\dfrac{2\left({\dfrac{1}{\sin^2 \alpha}}\right)-3\cot\alpha \left({\dfrac{1}{\sin^2 \alpha}}\right)}{5+\cot^3\alpha}\\
			&=\dfrac{-3\cot^3\alpha+2\cot^2\alpha-3\cot \alpha+2  }{5+\cot^3\alpha}\\
			&=-\dfrac{35}{16}. 
		\end{align*}
	} 
\end{bt}

% nguồn: Ok_TuanLuan_Toan11_Lan2
\begin{ex}%[0D6K2-2]%[ma:9fa248d9]
	Cho $\cot\alpha =3$. Tính giá trị biểu thức  $M=\dfrac{2\sin \alpha-3\cos\alpha}{5\sin^3 \alpha+\cos^3 \alpha}$.
	\loigiai{ 
		Ta  có
		\begin{align*}
			M&=\dfrac{2\sin  \alpha-3\cos\alpha}{5\sin^3 \alpha+\cos^3 \alpha}\\
			&=\dfrac{2\left({\dfrac{1}{\sin^2 \alpha}}\right)-3\cot\alpha  \left({\dfrac{1}{\sin^2 \alpha}}\right)}{5+\cot^3\alpha}\\
			&=\dfrac{-3\cot^3\alpha+2\cot^2\alpha-3\cot \alpha+2 }{5+\cot^3\alpha}\\ 
			&=-\dfrac{35}{16}. 
		\end{align*}
	 }
\end{ex}
